\documentclass[11pt, a4paper]{article}

\usepackage[a4paper, margin=2cm]{geometry}
\usepackage[spanish]{babel}
\usepackage{lmodern}
\usepackage[scaled=0.95]{helvet}
\usepackage{courier}

\usepackage{amsmath}
\usepackage{booktabs}
\usepackage{xcolor}
\usepackage{graphicx}
\usepackage{listings}
\usepackage{tcolorbox}
\tcbuselibrary{skins, breakable, listings}
\usepackage{hyperref}

% ==========================================
% PALETA DE COLORES DE ALTO CONTRASTE Y VISIBILIDAD
% ==========================================
\definecolor{brandNavy}{RGB}{15, 23, 42}       % Deep Navy Dark
\definecolor{brandBlue}{RGB}{37, 99, 235}      % Vibrant Blue Accent
\definecolor{brandCyan}{RGB}{14, 165, 233}     % Cyan Light Accent
\definecolor{codeBg}{RGB}{248, 250, 252}        % Canvas Soft Light Gray
\definecolor{codeFrame}{RGB}{203, 213, 225}     % Border Slate Gray
\definecolor{codeComment}{RGB}{100, 116, 139}   % Muted Comment Text
\definecolor{codeKeyword}{RGB}{3, 105, 161}     % Deep Teal-Blue Keyword
\definecolor{codeString}{RGB}{180, 83, 9}       % Warm Amber String
\definecolor{termBg}{RGB}{15, 23, 42}          % Dark Terminal Canvas
\definecolor{termBorder}{RGB}{30, 41, 59}      % Terminal Border
\definecolor{termFg}{RGB}{241, 245, 249}        % Terminal Text Light
\definecolor{boxTitleBg}{RGB}{30, 58, 138}     % Box Header Dark Blue

% ==========================================
% CONFIGURACIÓN DE LISTINGS Y TCOLORBOX
% ==========================================
\lstdefinestyle{codigo}{
  language=C++,
  basicstyle=\small\ttfamily\color{brandNavy},
  keywordstyle=\bfseries\color{codeKeyword},
  stringstyle=\color{codeString},
  commentstyle=\itshape\color{codeComment},
  numbers=left,
  numberstyle=\tiny\bfseries\color{codeComment},
  stepnumber=1,
  numbersep=10pt,
  showstringspaces=false,
  breaklines=true,
  tabsize=2
}

% Caja estilizada para código C / C++
\newtcblisting{codebox}[1][]{
  toptitle=2mm,
  bottomtitle=2mm,
  colback=codeBg,
  colframe=brandBlue,
  coltitle=white,
  fonttitle=\bfseries\small\sffamily,
  enhanced,
  attach boxed title to top left={xshift=5mm, yshift=-3mm},
  boxed title style={colback=brandBlue, arc=2mm},
  arc=3mm,
  boxrule=1.2pt,
  breakable,
  listing only,
  listing engine=listings,
  listing options={style=codigo},
  #1
}

% Caja estilizada para simulación de salida de terminal
\newtcblisting{terminalout}[1][]{
  toptitle=2mm,
  bottomtitle=2mm,
  colback=termBg,
  colframe=termBorder,
  colupper=termFg,
  fontupper=\small\ttfamily,
  fonttitle=\bfseries\small\sffamily,
  title={> COMPILACIÓN Y EJECUCIÓN EN CONSOLA},
  coltitle=white,
  enhanced,
  attach boxed title to top left={xshift=5mm, yshift=-3mm},
  boxed title style={colback=boxTitleBg, arc=2mm},
  arc=3mm,
  boxrule=1.2pt,
  breakable,
  listing only,
  listing options={
    basicstyle=\small\ttfamily\color{termFg},
    breaklines=true,
    showstringspaces=false
  },
  #1
}

\hypersetup{
    colorlinks=true,
    linkcolor=brandBlue,
    urlcolor=brandBlue,
    pdftitle={CAPÍTULO I - Solución de ejercicios de programación},
    pdfauthor={Fred Torres-Cruz}
}

\begin{document}

% ==========================================
% PORTADA / COVER PAGE CON DOS LOGOS
% ==========================================
\begin{titlepage}
    \centering
    
    % Encabezado con los dos logos alineados en los extremos (con verificación de existencia)
    \begin{minipage}{0.42\textwidth}
        \begin{center}
            \IfFileExists{logo-una-puno.png}{%
                \includegraphics[height=3.2cm, keepaspectratio]{logo-una-puno.png}%
            }{%
                \begin{tcolorbox}[colback=brandNavy, colframe=brandBlue, width=3.2cm, height=3.2cm, arc=2mm, halign=center, valign=center, boxrule=1pt]
                    {\small \bfseries \color{white} LOGO\\UNA-PUNO}
                \end{tcolorbox}%
            }
        \end{center}
    \end{minipage}
    \hfill
    \begin{minipage}{0.42\textwidth}
        \begin{center}
            \IfFileExists{logofinesi.png}{%
                \includegraphics[height=3.2cm, keepaspectratio]{logofinesi.png}%
            }{%
                \begin{tcolorbox}[colback=brandNavy, colframe=brandBlue, width=3.2cm, height=3.2cm, arc=2mm, halign=center, valign=center, boxrule=1pt]
                    {\small \bfseries \color{white} LOGO\\FINESI}
                \end{tcolorbox}%
            }
        \end{center}
    \end{minipage}

    \vspace{1.5cm}

    {\Large \bfseries UNIVERSIDAD NACIONAL DEL ALTIPLANO \par}
    {\large FACULTAD DE INGENIERÍA ESTADÍSTICA E INFORMÁTICA \par}

    \vspace{1.5cm}

    \begin{tcolorbox}[colback=brandNavy, colframe=brandBlue, arc=4mm, width=\textwidth, halign=center, top=8mm, bottom=8mm]
        {\huge \bfseries \color{white} CAPÍTULO I \par}
        \vspace{0.3cm}
        {\Large \color{brandCyan} Solución de Ejercicios de Programación \par}
    \end{tcolorbox}
    
    \vspace{1.2cm}
    
    {\huge \bfseries \color{brandNavy} ESTRUCTURA DE DATOS -- 2026 \par}
    \vspace{0.4cm}
    \rule{0.8\textwidth}{1.5pt}
    \vspace{0.8cm}
    
    {\Large \textbf{ Docente:} \color{brandBlue} Fred Torres-Cruz \par}
    {\Large \textbf{ Estudiante:} \color{brandBlue} Blays Kevin Mamani Quispe \par}
    {\Large \textbf{ Link de Repositorio GitHub:} \color{brandBlue} https://github.com/blayskevinmamaniquispe-lab/Estructura-de-Datos/tree/main/Ejerccicios%20-%20Estructuras%20b%C3%A1sicas. \par}
    \vfill
    
    \begin{tcolorbox}[colback=codeBg, colframe=brandBlue, arc=3mm, width=0.9\textwidth, halign=left, boxrule=1pt]
        \textbf{\large Resumen del Capítulo:}
        \begin{itemize}\setlength{\itemsep}{2pt}
            \item Ejercicios resueltos en lenguaje C y C++.
            \item Estructuras secuenciales, condicionales, ciclos y arreglos.
            \item Modularización mediante funciones y estimación de complejidad algorítmica.
        \end{itemize}
    \end{tcolorbox}
    
    \vfill
    
    {\large \bfseries \color{codeComment} Puno, Perú \par}
    {\large \color{codeComment} 2026 \par}
\end{titlepage}

\newpage
\tableofcontents
\newpage

% ==========================================
% DESARROLLO DE EJERCICIOS
% ==========================================

\section*{Ejercicio 1. Costo total del procesamiento}
\addcontentsline{toc}{section}{Ejercicio 1. Costo total del procesamiento}
\textbf{Procedimiento:} Se leen las horas trabajadas y el costo por hora. El costo total se obtiene multiplicando ambos valores.

\begin{codebox}[title={CÓDIGO FUENTE EN C: Ejercicio 1}]
#include <stdio.h>

int main(void) {
    float horas, tarifa, costo;

    printf("Ingrese las horas trabajadas: ");
    scanf("%f", &horas);

    printf("Ingrese el costo por hora: ");
    scanf("%f", &tarifa);

    costo = horas * tarifa;

    printf("\n--- RESUMEN ---\n");
    printf("Horas trabajadas: %.2f\n", horas);
    printf("Tarifa por hora: S/ %.2f\n", tarifa);
    printf("Costo total: S/ %.2f\n", costo);

    return 0;
}
\end{codebox}

\begin{terminalout}
$ gcc ejercicio1.c -o ejercicio1
$ ./ejercicio1
Ingrese las horas trabajadas: 42.5
Ingrese el costo por hora: 25.0

--- RESUMEN ---
Horas trabajadas: 42.50
Tarifa por hora: S/ 25.00
Costo total: S/ 1062.50
\end{terminalout}

\newpage

\section*{Ejercicio 2. Registro de una observación}
\addcontentsline{toc}{section}{Ejercicio 2. Registro de una observación}
\textbf{Procedimiento:} Se utilizan tipos de datos apropiados: entero para el identificador y la edad, float para el promedio, char para la categoría y un entero para representar el estado de validez.

\begin{codebox}[title={CÓDIGO FUENTE EN C: Ejercicio 2}]
#include <stdio.h>

int main(void) {
    int identificador, edad, valido;
    float promedio;
    char categoria;

    printf("Identificador: ");
    scanf("%d", &identificador);

    printf("Edad: ");
    scanf("%d", &edad);

    printf("Valor promedio: ");
    scanf("%f", &promedio);

    printf("Categoria (una letra): ");
    scanf(" %c", &categoria);

    printf("Registro valido? (1=Si, 0=No): ");
    scanf("%d", &valido);

    printf("\n--- DATOS DE LA OBSERVACION ---\n");
    printf("Identificador: %d\n", identificador);
    printf("Edad: %d\n", edad);
    printf("Promedio: %.2f\n", promedio);
    printf("Categoria: %c\n", categoria);
    printf("Estado: %s\n", valido == 1 ? "Valido" : "No valido");

    return 0;
}
\end{codebox}

\begin{terminalout}
$ gcc ejercicio2.c -o ejercicio2
$ ./ejercicio2
Identificador: 101
Edad: 30
Valor promedio: 56.45
Categoria (una letra): B
Registro valido? (1=Si, 0=No): 1

--- DATOS DE LA OBSERVACION ---
Identificador: 101
Edad: 30
Promedio: 56.45
Categoria: B
Estado: Valido
\end{terminalout}

\newpage

\section*{Ejercicio 3. Conversión de decimal a entero}
\addcontentsline{toc}{section}{Ejercicio 3. Conversión de decimal a entero}
\textbf{Procedimiento:} La conversión explícita mediante \texttt{(int)} elimina la parte decimal, sin redondear. La pérdida se calcula restando el entero convertido al valor original.

\begin{codebox}[title={CÓDIGO FUENTE EN C: Ejercicio 3}]
#include <stdio.h>

int main(void) {
    float medicion, perdida;
    int valorEntero;

    printf("Ingrese la medicion decimal: ");
    scanf("%f", &medicion);

    valorEntero = (int)medicion;
    perdida = medicion - valorEntero;

    printf("Valor original: %.2f\n", medicion);
    printf("Valor entero: %d\n", valorEntero);
    printf("Parte decimal perdida: %.2f\n", perdida);

    return 0;
}
\end{codebox}

\begin{terminalout}
$ gcc ejercicio3.c -o ejercicio3
$ ./ejercicio3
Ingrese la medicion decimal: 18.9

Valor original: 18.90
Valor entero: 18
Parte decimal perdida: 0.90
\end{terminalout}

\newpage

\section*{Ejercicio 4. Operaciones aritméticas}
\addcontentsline{toc}{section}{Ejercicio 4. Operaciones aritméticas}
\textbf{Procedimiento:} Se leen dos enteros y se calculan suma, diferencia, producto, división entera y residuo. Se comprueba que el segundo número no sea cero antes de dividir.

\begin{codebox}[title={CÓDIGO FUENTE EN C: Ejercicio 4}]
#include <stdio.h>

int main(void) {
    int a, b;

    printf("Ingrese el primer entero: ");
    scanf("%d", &a);

    printf("Ingrese el segundo entero: ");
    scanf("%d", &b);

    printf("\nSuma: %d\n", a + b);
    printf("Diferencia: %d\n", a - b);
    printf("Producto: %d\n", a * b);

    if (b != 0) {
        printf("Division entera: %d\n", a / b);
        printf("Residuo: %d\n", a % b);
    } else {
        printf("No se puede dividir entre cero.\n");
    }

    return 0;
}
\end{codebox}

\begin{terminalout}
$ gcc ejercicio4.c -o ejercicio4
$ ./ejercicio4
Ingrese el primer entero: 150
Ingrese el segundo entero: 12

Suma: 162
Diferencia: 138
Producto: 1800
Division entera: 12
Residuo: 6
\end{terminalout}

\newpage

\section*{Ejercicio 5. Distribución de registros entre nodos}
\addcontentsline{toc}{section}{Ejercicio 5. Distribución de registros entre nodos}
\textbf{Procedimiento:} La división entera determina los registros asignados a cada nodo y el operador módulo determina los registros sobrantes.

\begin{codebox}[title={CÓDIGO FUENTE EN C: Ejercicio 5}]
#include <stdio.h>

int main(void) {
    int registros, nodos;

    printf("Numero total de registros: ");
    scanf("%d", &registros);

    printf("Numero de nodos: ");
    scanf("%d", &nodos);

    if (registros >= 0 && nodos > 0) {
        printf("Registros por nodo: %d\n", registros / nodos);
        printf("Registros sin distribuir: %d\n", registros % nodos);
    } else {
        printf("Ingrese registros no negativos y nodos mayores que cero.\n");
    }

    return 0;
}
\end{codebox}

\begin{terminalout}
$ gcc ejercicio5.c -o ejercicio5
$ ./ejercicio5
Numero total de registros: 1053
Numero de nodos: 8

Registros por nodo: 131
Registros sin distribuir: 5
\end{terminalout}

\newpage

\section*{Ejercicio 6. Operadores de incremento y asignación}
\addcontentsline{toc}{section}{Ejercicio 6. Operadores de incremento y asignación}
\textbf{Procedimiento:} La variable inicia en 4, luego aumenta en 3 y finalmente se multiplica por 2.

\begin{codebox}[title={CÓDIGO FUENTE EN C: Ejercicio 6}]
#include <stdio.h>

int main(void) {
    int registros = 4;

    printf("Valor inicial: %d\n", registros);

    registros += 3;
    printf("Despues de incrementar en 3: %d\n", registros);

    registros *= 2;
    printf("Despues de duplicar: %d\n", registros);

    return 0;
}
\end{codebox}

\begin{terminalout}
$ gcc ejercicio6.c -o ejercicio6
$ ./ejercicio6
Valor inicial: 4
Despues de incrementar en 3: 7
Despues de duplicar: 14
\end{terminalout}

\newpage

\section*{Ejercicio 7. Validación de un conjunto de datos}
\addcontentsline{toc}{section}{Ejercicio 7. Validación de un conjunto de datos}
\textbf{Procedimiento:} El conjunto se acepta únicamente si tiene al menos 100 observaciones válidas y un porcentaje de datos completos igual o superior al 70\%. Ambas condiciones deben cumplirse.

\begin{codebox}[title={CÓDIGO FUENTE EN C: Ejercicio 7}]
#include <stdio.h>

int main(void) {
    int observaciones;
    float porcentaje;

    printf("Numero de observaciones validas: ");
    scanf("%d", &observaciones);

    printf("Porcentaje de datos completos: ");
    scanf("%f", &porcentaje);

    if (observaciones >= 100 && porcentaje >= 70.0 && porcentaje <= 100.0) {
        printf("El conjunto de datos puede ser analizado.\n");
    } else {
        printf("El conjunto de datos no cumple las condiciones.\n");
    }

    return 0;
}
\end{codebox}

\begin{terminalout}
$ gcc ejercicio7.c -o ejercicio7
$ ./ejercicio7
Numero de observaciones validas: 125
Porcentaje de datos completos: 82.5

El conjunto de datos puede ser analizado.
\end{terminalout}

\newpage

\section*{Ejercicio 8. Validación de una observación}
\addcontentsline{toc}{section}{Ejercicio 8. Validación de una observación}
\textbf{Procedimiento:} Se utiliza el operador lógico AND (\texttt{\&\&}). La observación se incorpora si la edad es como mínimo 18 y el registro está validado.

\begin{codebox}[title={CÓDIGO FUENTE EN C: Ejercicio 8}]
#include <stdio.h>

int main(void) {
    int edad, validado;

    printf("Ingrese la edad: ");
    scanf("%d", &edad);

    printf("Registro validado? (1=Si, 0=No): ");
    scanf("%d", &validado);

    if (edad >= 18 && validado == 1) {
        printf("La observacion puede incorporarse.\n");
    } else {
        printf("La observacion no puede incorporarse.\n");
    }

    return 0;
}
\end{codebox}

\begin{terminalout}
$ gcc ejercicio8.c -o ejercicio8
$ ./ejercicio8
Ingrese la edad: 22
Registro validado? (1=Si, 0=No): 1

La observacion puede incorporarse.
\end{terminalout}

\newpage

\section*{Ejercicio 9. Clasificación de temperatura}
\addcontentsline{toc}{section}{Ejercicio 9. Clasificación de temperatura}
\textbf{Procedimiento:} Se utilizan condiciones encadenadas para clasificar la temperatura en tres categorías.

\begin{codebox}[title={CÓDIGO FUENTE EN C: Ejercicio 9}]
#include <stdio.h>

int main(void) {
    float temperatura;

    printf("Ingrese la temperatura en grados Celsius: ");
    scanf("%f", &temperatura);

    if (temperatura < 0) {
        printf("Clasificacion: Congelacion\n");
    } else if (temperatura <= 20) {
        printf("Clasificacion: Frio\n");
    } else {
        printf("Clasificacion: Templado\n");
    }

    return 0;
}
\end{codebox}

\begin{terminalout}
$ gcc ejercicio9.c -o ejercicio9
$ ./ejercicio9
Ingrese la temperatura en grados Celsius: 14.5
Clasificacion: Frio
\end{terminalout}

\newpage

\section*{Ejercicio 10. Media y observaciones superiores a ella (C++)}
\addcontentsline{toc}{section}{Ejercicio 10. Media y observaciones superiores a ella (C++)}
\textbf{Procedimiento:} Se almacenan diez mediciones en un arreglo y se calcula su suma. La media se obtiene dividiendo la suma entre diez. Después se cuentan los valores mayores que la media.

\begin{codebox}[title={CÓDIGO FUENTE EN C++: Ejercicio 10}]
#include <iostream>
using namespace std;

int main() {
    double datos[10], suma = 0, media;
    int contador = 0;

    for (int i = 0; i < 10; i++) {
        cout << "Ingrese la medicion " << i + 1 << ": ";
        cin >> datos[i];
        suma += datos[i];
    }

    media = suma / 10;

    for (int i = 0; i < 10; i++) {
        if (datos[i] > media) {
            contador++;
        }
    }

    cout << "Media aritmetica: " << media << endl;
    cout << "Observaciones por encima de la media: " << contador << endl;

    return 0;
}
\end{codebox}

\begin{terminalout}
$ g++ ejercicio10.cpp -o ejercicio10
$ ./ejercicio10
Ingrese la medicion 1: 12.5
Ingrese la medicion 2: 18.0
Ingrese la medicion 3: 9.2
... [ingreso de mediciones] ...
Media aritmetica: 14.82
Observaciones por encima de la media: 4
\end{terminalout}

\newpage

\section*{Ejercicio 11. Validación de contraseña (C++)}
\addcontentsline{toc}{section}{Ejercicio 11. Validación de contraseña (C++)}
\textbf{Procedimiento:} El ciclo \texttt{while} solicita la clave hasta que sea correcta. Se incrementa el contador cada vez que se realiza un intento. Para este ejemplo, la clave correcta se establece en 1234.

\begin{codebox}[title={CÓDIGO FUENTE EN C++: Ejercicio 11}]
#include <iostream>
using namespace std;

int main() {
    const int CLAVE_CORRECTA = 1234;
    int clave, intentos = 0;

    cout << "Ingrese la clave numerica: ";
    cin >> clave;
    intentos++;

    while (clave != CLAVE_CORRECTA) {
        cout << "Clave incorrecta. Intente nuevamente.\n";
        cout << "Ingrese la clave numerica: ";
        cin >> clave;
        intentos++;
    }

    cout << "Acceso permitido.\n";
    cout << "Numero de intentos: " << intentos << endl;

    return 0;
}
\end{codebox}

\begin{terminalout}
$ g++ ejercicio11.cpp -o ejercicio11
$ ./ejercicio11
Ingrese la clave numerica: 1111
Clave incorrecta. Intente nuevamente.
Ingrese la clave numerica: 1234
Acceso permitido.
Numero de intentos: 2
\end{terminalout}

\newpage

\section*{Ejercicio 12. Procesamiento con break y continue (C++)}
\addcontentsline{toc}{section}{Ejercicio 12. Procesamiento con break y continue (C++)}
\textbf{Procedimiento:} Se leen como máximo diez valores. Los negativos se omiten con \texttt{continue}; el valor 999 detiene el ciclo con \texttt{break}.

\begin{codebox}[title={CÓDIGO FUENTE EN C++: Ejercicio 12}]
#include <iostream>
using namespace std;

int main() {
    double valor;
    int validos = 0;

    for (int i = 0; i < 10; i++) {
        cout << "Ingrese la medicion " << i + 1 << " (999 para terminar): ";
        cin >> valor;

        if (valor == 999) {
            break;
        }

        if (valor < 0) {
            cout << "Valor negativo: se omite.\n";
            continue;
        }

        validos++;
    }

    cout << "Valores validos procesados: " << validos << endl;

    return 0;
}
\end{codebox}

\begin{terminalout}
$ g++ ejercicio12.cpp -o ejercicio12
$ ./ejercicio12
Ingrese la medicion 1 (999 para terminar): 15.4
Ingrese la medicion 2 (999 para terminar): -3.2
Valor negativo: se omite.
Ingrese la medicion 3 (999 para terminar): 999
Valores validos procesados: 1
\end{terminalout}

\newpage

\section*{Ejercicio 13. Análisis de edades mediante un arreglo}
\addcontentsline{toc}{section}{Ejercicio 13. Análisis de edades mediante un arreglo}
\textbf{Procedimiento:} Se almacenan diez edades. Durante el ingreso se determina la suma, la edad mínima y la máxima. Después se calcula la media y se cuentan las edades superiores a ella.

\begin{codebox}[title={CÓDIGO FUENTE EN C: Ejercicio 13}]
#include <stdio.h>

int main(void) {
    int edades[10], suma = 0;
    int minima, maxima, superiores = 0;
    float media;

    for (int i = 0; i < 10; i++) {
        printf("Edad del participante %d: ", i + 1);
        scanf("%d", &edades[i]);

        if (i == 0) {
            minima = maxima = edades[i];
        } else {
            if (edades[i] < minima) minima = edades[i];
            if (edades[i] > maxima) maxima = edades[i];
        }

        suma += edades[i];
    }

    media = (float)suma / 10;

    for (int i = 0; i < 10; i++) {
        if (edades[i] > media) {
            superiores++;
        }
    }

    printf("\nEdad minima: %d\n", minima);
    printf("Edad maxima: %d\n", maxima);
    printf("Media de edades: %.2f\n", media);
    printf("Participantes por encima de la media: %d\n", superiores);

    return 0;
}
\end{codebox}

\begin{terminalout}
$ gcc ejercicio13.c -o ejercicio13
$ ./ejercicio13
Edad minima: 19
Edad maxima: 45
Media de edades: 28.40
Participantes por encima de la media: 4
\end{terminalout}

\newpage

\section*{Ejercicio 14. Comparación de complejidad \texorpdfstring{$O(n)$}{O(n)} y \texorpdfstring{$O(n^2)$}{O(n^2)}}
\addcontentsline{toc}{section}{Ejercicio 14. Comparación de complejidad O(n) y O(n2)}
\textbf{Procedimiento:} El primer procedimiento ejecuta una operación por cada elemento, por lo que realiza $n$ operaciones. El segundo ejecuta dos ciclos anidados de $n$ iteraciones cada uno, realizando $n^2$ operaciones.

\begin{codebox}[title={CÓDIGO FUENTE EN C: Ejercicio 14}]
#include <stdio.h>

int main(void) {
    int n;
    long long operacionesLineales = 0;
    long long operacionesCuadraticas = 0;

    printf("Ingrese el numero de elementos: ");
    scanf("%d", &n);

    if (n <= 0) {
        printf("n debe ser un entero positivo.\n");
        return 1;
    }

    /* Procedimiento 1: un recorrido */
    for (int i = 0; i < n; i++) {
        operacionesLineales++;
    }

    /* Procedimiento 2: dos ciclos anidados */
    for (int i = 0; i < n; i++) {
        for (int j = 0; j < n; j++) {
            operacionesCuadraticas++;
        }
    }

    printf("Operaciones del primer procedimiento: %lld\n", operacionesLineales);
    printf("Complejidad aproximada: O(n)\n");

    printf("Operaciones del segundo procedimiento: %lld\n", operacionesCuadraticas);
    printf("Complejidad aproximada: O(n^2)\n");

    return 0;
}
\end{codebox}

\begin{terminalout}
$ gcc ejercicio14.c -o ejercicio14
$ ./ejercicio14
Ingrese el numero de elementos: 500
Operaciones del primer procedimiento: 500
Complejidad aproximada: O(n)
Operaciones del segundo procedimiento: 250000
Complejidad aproximada: O(n^2)
\end{terminalout}

\newpage

\section*{Ejercicio 15. Funciones para promedio y conteo}
\addcontentsline{toc}{section}{Ejercicio 15. Funciones para promedio y conteo}
\textbf{Procedimiento:} La función \texttt{calcular\_promedio} recibe tres mediciones y retorna su media. La función \texttt{contar\_superiores} recibe las mediciones y el promedio, retornando cuántas superan dicho promedio.

\begin{codebox}[title={CÓDIGO FUENTE EN C: Ejercicio 15}]
#include <stdio.h>

float calcular_promedio(float a, float b, float c) {
    return (a + b + c) / 3.0f;
}

int contar_superiores(float a, float b, float c, float promedio) {
    int contador = 0;

    if (a > promedio) contador++;
    if (b > promedio) contador++;
    if (c > promedio) contador++;

    return contador;
}

int main(void) {
    float m1, m2, m3, promedio;
    int cantidad;

    printf("Ingrese la primera medicion: ");
    scanf("%f", &m1);

    printf("Ingrese la segunda medicion: ");
    scanf("%f", &m2);

    printf("Ingrese la tercera medicion: ");
    scanf("%f", &m3);

    promedio = calcular_promedio(m1, m2, m3);
    cantidad = contar_superiores(m1, m2, m3, promedio);

    printf("Promedio: %.2f\n", promedio);
    printf("Mediciones por encima del promedio: %d\n", cantidad);

    return 0;
}
\end{codebox}

\begin{terminalout}
$ gcc ejercicio15.c -o ejercicio15
$ ./ejercicio15
Ingrese la primera medicion: 14.2
Ingrese la segunda medicion: 18.6
Ingrese la tercera medicion: 11.0
Promedio: 14.60
Mediciones por encima del promedio: 1
\end{terminalout}

\newpage

% CONCLUSION
\section*{Conclusión}
\addcontentsline{toc}{section}{Conclusión}

Los ejercicios desarrollan conceptos fundamentales de programación: variables y tipos de datos, conversiones, operadores aritméticos y lógicos, estructuras condicionales, ciclos, arreglos, funciones y análisis básico de complejidad algorítmica.

Cada programa puede compilarse de manera independiente en Overleaf o cualquier compilador de C/C++.

\end{document}